\section{Introduction}
Surface electromyography (sEMG) provides a non-invasive signal for gesture recognition and myoelectric control, but the recording interface must balance spatial information against the burden of additional electrodes \citep{scheme2011electromyogram,farina2014extraction,vujaklija2016new}. Fewer recording sites could simplify an interface. Whether that reduction preserves recognition for an unseen participant is an empirical question, and offline recognition alone does not establish a usable control interface.

Channel selection has been studied with feature filters, wrappers and spatial criteria \citep{phinyomark2012feature,stango2015spatial,qu2021reduce}. SHAP-guided EMG channel and feature optimization is also established \citep{rani2025shap}; related work uses Shapley explanations for similar sEMG gestures \citep{wang2024shap} and explainable feature selection with fusion \citep{gehlot2024fusion}. These contributions make attribution-based reduction a relevant comparator rather than a new algorithmic premise. Their performance estimates cannot be directly compared with this study across different tasks, sensors and validation protocols.

For a system intended to work across people, selection must exclude the eventual test participant. Using that participant to standardize features, estimate importance or choose a subset can produce optimistic estimates. Participant-disjoint inner evaluation also avoids measuring attribution solely on observations used to fit the attributing model. This motivates an outer leave-one-subject-out (LOSO) evaluation in which each entire selection procedure is reconstructed from the outer training participants.

Attribution magnitude describes predictive reliance of a fitted model, not a direct measurement of muscle activation. TreeSHAP distributes a model's prediction over its input features \citep{lundberg2017unified,lundberg2020local}; aggregation over an electrode's features offers a ranking, but does not turn it into a coalition-level Shapley value for that electrode. Correlated features and variation across participants can affect this ranking. Attribution stability has already motivated resampling and regularization methods \citep{visani2022statistical,xiang2026roshap}; a stability penalty therefore requires validation rather than a novelty claim.

Electrode count requires an explicit tolerance for performance loss. Non-inferiority asks whether a reduced system remains within such a tolerance, whereas a non-significant superiority test does not establish preservation \citep{schuirmann1987comparison,walker2011understanding}. A mean tolerance is also different from preserving every participant or gesture. Reporting only an aggregate curve could hide losses that matter to an intended application.

We evaluate cross-fitted electrode reduction on GRABMyo and separately examine NinaPro DB5. The primary question is the minimum count supported by an internally documented, predefined engineering margin. Secondary questions concern utility beyond random and spatial selection, stability regularization, and deletion with retraining. New exploratory checks use a geometric subset under the primary training protocol, paired participant and gesture predictions, and shared-resource SHAP and grouped permutation rankings. Sensitivities and split-derived layouts delimit the claim: the LOSO experiment evaluates a selection procedure, not a prospectively validated common physical layout.

\begin{figure}[htbp]\centering
\includegraphics[width=\textwidth]{F:/grabmyo_data/manuscript_revision/20261014/figures/vector/fig1_montage.pdf}
\caption{Recording topology and task. The source montage places F1--F8 in the proximal ring and F9--F16 in the distal ring, separated by 2 cm. Corresponding circumferential positions are aligned. The forearm outline is schematic; recording sites were not registered to individual muscles. The unwrapped map preserves electrode numbers without implying a verified absolute anatomical angle.}
\label{fig:montage}\end{figure}
